\documentclass[professionalfonts,aspectratio=169]{beamer}
\usepackage{newtxtext,newtxmath}
\usepackage[utf8]{inputenc}
\usepackage[brazil]{babel}
\usepackage[T1]{fontenc} 
\usepackage{graphicx} 
\usepackage{bibentry} 
\usepackage{txfonts}			% Fontes virtuais
\setbeamertemplate{caption}[numbered]
\usepackage{caption}
\usepackage[alf]{abntex2cite}		% Citações padrão ABNT


% Remove os símbolos de navegação
\usenavigationsymbolstemplate{}

% Theme choice:
\usetheme{Boadilla}

% Cria a cor de padrão da marca (azul, laranja, branco e preto)
\definecolor{UniBlue}{RGB}{32,64,154}
\definecolor{UniOrange}{cmyk}{0,50,100,0}

% Define as cores dos parametros
\setbeamercolor{title}{fg=UniBlue}
\setbeamercolor{author}{fg=UniBlue}
\setbeamercolor{institute}{fg=UniBlue}
\setbeamercolor{curso}{fg=UniBlue}
\setbeamercolor{frametitle}{fg=UniBlue}
\setbeamercolor{structure}{fg=UniBlue}
\setbeamercolor{normal text}{fg=black}
\usebeamercolor[fg]{normal text}
\setbeamercolor{enumerate}{fg=black}
\setbeamercolor{itemize}{fg=black}
\setbeamercolor{section in toc}{fg=black}
\setbeamercolor{subsection in toc}{fg=black}

% Define os tamanhos dos parametros
% size*={<font size>}{<baseline skip>}
\setbeamerfont{title}{size=\fontsize{12pt}{8pt}}
\setbeamerfont{author}{size=\fontsize{16pt}{8pt}}
\setbeamerfont{institute}{size=\fontsize{16pt}{0pt}}
\setbeamerfont{curso}{size=\fontsize{14pt}{0pt}}

\addtobeamertemplate{institute}{\vspace*{20\baselineskip}}{}


\setbeamerfont{date}{size=\fontsize{8pt}{0pt}}

\setbeamerfont{frametitle}{size=\fontsize{16pt}{8pt}}
\addtobeamertemplate{frametitle}{\vspace*{1\baselineskip}}{}

\setbeamerfont{normal text}{size=\fontsize{10pt}{8pt}}
\setbeamerfont{enumerate}{size=\fontsize{10pt}{8pt}}
\setbeamerfont{itemize}{size=\fontsize{10pt}{8pt}}
\setbeamerfont{section in toc}{size=\fontsize{10pt}{8pt}}
\setbeamerfont{subsection in toc}{size=\fontsize{8pt}{8pt}}

%%%%%%%%%%%%%%%%%%%%%%%%%%%%%%%%%%%
\setbeamertemplate{title page}
{
    \vbox{}
    \vfill
    \begingroup
        \centering
        %-----------------institute
        \begin{beamercolorbox}[sep=8pt,right]{institute}
            \usebeamerfont{institute}\insertinstitute
        \end{beamercolorbox}
        %-----------------title
        \begin{beamercolorbox}[sep=8pt,right]{title}
            \usebeamerfont{title}\inserttitle\par%
        \end{beamercolorbox}%
        %-----------------Autor
       % \begin{beamercolorbox}[sep=8pt,right]{author}
       %     \usebeamerfont{author}\insertauthor
       % \end{beamercolorbox}
    \endgroup
    \vfill
}

%%%%%%%%%%%%%%%%%%%%%%%%%%%%%%%%%%%
% global background must be put in the preamble
\setbeamertemplate{background canvas}
{%
    \includegraphics[
    page=1,
    %trim = 20mm 0mm 0mm -20mm, % use esse corte em caso de aspecto 4:3
    width=\paperwidth, height=\paperheight]{logos/Unisinos.pdf}%
}
%%%%%%%%%%%%%%%%%%%%%%%%%%%%%%%%%%%


% Dados
\title{TÍTULO DO TRABALHO}
\author{Fulano de Tal}
%\curso{Engenharia Eletrônica}% \\[5mm] ORIENTADOR: PROF. DR. FULANO DE TAL}
%\institute{Universidade do Vale do Rio dos Sinos - UNISINOS}
\institute{NOME DO CURSO}
%\date{\\[10mm] 12 DE AGOSTO DE 2017}

\newcommand*\oldmacro{}%
\let\oldmacro\insertshorttitle%
\renewcommand*\insertshorttitle{%
  \oldmacro\hfill%
  \insertframenumber\,/\,\inserttotalframenumber}
\begin{document}
\def\beamer@andinst{\\[5em]}
%%%%%%%%%%%%%%%%%%%%%%%%%%%%%%%%%%%
\frame{\titlepage}
%%%%%%%%%%%%%%%%%%%%%%%%%%%%%%%%%%%

\setbeamertemplate{background canvas}{
    \includegraphics[
    page=2,
    %trim = 20mm 0mm 5mm -20mm, % use esse corte em caso de aspecto 4:3
    width=\paperwidth, height=\paperheight,
    ]{logos/Unisinos.pdf}%
}
%%%%%%%%%%%%%%%%%%%%%%%%%%%%%%%%%%%
\section*{Sumário}
%%%%%%%%%%%%%%%%%%%%%%%%%%%%%%%%%%%
\begin{frame}{Sumário}
	\tableofcontents[]
\end{frame}

%%%%%%%%%%%%%%%%%%%%%%%%%%%%%%%%%%%
\section{Introdução}
%%%%%%%%%%%%%%%%%%%%%%%%%%%%%%%%%%%
\begin{frame}{Introdução} 
A introdução do trabalho vai neste slide
\end{frame}

%%%%%%%%%%%%%%%%%%%%%%%%%%%%%%%%%%%
\section{Math equations}
%%%%%%%%%%%%%%%%%%%%%%%%%%%%%%%%%%%
\begin{frame}{Math equations}
\begin{block}{Math equations}
The Pythagorean theorem is a fundamental relation in Euclidean geometry among the three sides of a right triangle describe by \citeonline{Nirmalasari2021}:
\begin{equation}
    a^{2}+b^{2}=c^{2}
\end{equation}  
\end{block}   
\end{frame}

%%%%%%%%%%%%%%%%%%%%%%%%%%%%%%%%%%%
\section{Proposta}
%%%%%%%%%%%%%%%%%%%%%%%%%%%%%%%%%%%
\begin{frame}{Proposta}
\begin{itemize}
\pause \item Tópico 1
\pause \item Mais um
\pause \item Outro
\end{itemize}
\end{frame}

%%%%%%%%%%%%%%%%%%%%%%%%%%%%%%%%%%%
\section{Experimentos}
%%%%%%%%%%%%%%%%%%%%%%%%%%%%%%%%%%%
\begin{frame}{Experimentos}
\begin{itemize}
\pause \item Resultados
\end{itemize}
\end{frame}

%%%%%%%%%%%%%%%%%%%%%%%%%%%%%%%%%%%
\section{Conclusão}
%%%%%%%%%%%%%%%%%%%%%%%%%%%%%%%%%%%
\begin{frame}{Conclusão}
\begin{itemize}
\pause \item Conclusão
\end{itemize}
\end{frame}

%%%%%%%%%%%%%%%%%%%%%%%%%%%%%%%%%%%
\section{Bibliografia}
%%%%%%%%%%%%%%%%%%%%%%%%%%%%%%%%%%%
\begin{frame}{Bibliografia}
    \bibliography{export}
\end{frame}

%%%%%%%%%%%%%%%%%%%%%%%%%%%%%%%%%%%
\section{Fim}
%%%%%%%%%%%%%%%%%%%%%%%%%%%%%%%%%%%
\setbeamertemplate{background canvas}{
    \includegraphics[
        page=3,
        %trim = 0mm 0mm 30mm -20mm, % use esse corte em caso de aspecto 4:3
        width=\paperwidth,height=\paperheight
        ]{logos/Unisinos.pdf}%
}
\begin{frame}{}
\end{frame}
%%%%%%%%%%%%%%%%%%%%%%%%%%%%%%%%%%%
\end{document}
